\documentclass[11pt]{article}
% Shared preamble for the printable mini lessons (1.4 onward).
%   \documentclass[11pt]{article}
%   \input{../preamble}
% Build the key with:  pdflatex -jobname=<name>_KEY "\def\iskey{}\input{<name>}"
% Figures shared with the web page live in figures/Unit_X/<name>.tex and are
% pulled in with \fig{Unit_X/<name>}; the standalone package drops their own
% preamble, so their macros come from figures/preamble.tex, loaded here.
% Shared preamble for web figures.  Each figure is a standalone .tex that
% \input's this file; scripts/make-figures.sh turns it into an SVG.
\usepackage[T1]{fontenc}
\usepackage{helvet}
\renewcommand{\familydefault}{\sfdefault}
\usepackage{tikz}
\usetikzlibrary{decorations.pathreplacing}
% axis numbers in the text font, not math italic
\pgfkeys{/pgf/number format/.cd,assume math mode=true,1000 sep={}}

% dot plot: \dotplot{min}{max}{value/count,...}{axis label}
\newcommand{\dotplot}[4]{%
\begin{tikzpicture}[x=0.9cm,y=0.33cm]
  \draw[thick] (#1-0.5,0) -- (#2+0.5,0);
  \foreach \t in {#1,...,#2} {\draw (\t,0.25) -- (\t,-0.25) node[below] {\t};}
  \foreach \v/\c in {#3} {\ifnum\c>0 \foreach \k in {1,...,\c} {\fill (\v,\k*1.15) circle (0.16cm);}\fi}
  \node at ({(#1+#2)/2},-2.3) {#4};
\end{tikzpicture}}

% shape sketch histogram (12 bins): \shapehist{heights}{label}
\newcommand{\shapehist}[2]{%
  \foreach \h [count=\i from 0] in {#1} \draw[fill=gray!30] (\i,0) rectangle (\i+1,\h);
  \draw (-0.3,0) -- (12.3,0);
  \node[font=\bfseries] at (6,-2.6) {#2};}
\newcommand{\meanmark}[1]{\fill (#1,0) -- ++(-0.3,-0.75) -- ++(0.6,0) -- cycle;}
\newcommand{\medmark}[2]{\draw[very thick,dashed] (#1,0) -- (#1,#2);}

% number line with minor ticks: \nline{min}{max}{label step}{axis label}
\newcommand{\nline}[4]{%
  \draw[thick] (#1-0.5,0) -- (#2+0.5,0);
  \foreach \t in {#1,...,#2} \draw (\t,0.12) -- (\t,-0.12);
  \pgfmathtruncatemacro{\nxt}{#1+#3}
  \foreach \t in {#1,\nxt,...,#2} \draw (\t,0.2) -- (\t,-0.2) node[below] {\t};
  \node at ({(#1+#2)/2},-0.95) {#4};}
% box plot above the line: \bplot{min}{Q1}{med}{Q3}{max}
\newcommand{\bplot}[5]{%
  \draw[thick] (#1,1.0) -- (#2,1.0);
  \draw[thick] (#4,1.0) -- (#5,1.0);
  \draw[thick] (#1,0.75) -- (#1,1.25);
  \draw[thick] (#5,0.75) -- (#5,1.25);
  \draw[thick,fill=gray!20] (#2,0.55) rectangle (#4,1.45);
  \draw[very thick] (#3,0.55) -- (#3,1.45);}

% ---- added for 1.4 onward ----
% axis on y=0: \axis{min}{max}{label step}{tick step}{axis label}
\newcommand{\axis}[5]{%
  \draw[thick] (#1,0) -- (#2,0);
  \pgfmathsetmacro{\axT}{#1+#4}\foreach \t in {#1,\axT,...,#2} \draw (\t,0.1) -- (\t,-0.1);
  \pgfmathsetmacro{\axL}{#1+#3}\foreach \t in {#1,\axL,...,#2}
    \draw (\t,0.18) -- (\t,-0.18) node[below] {\pgfmathprintnumber{\t}};
  \node at ({(#1+#2)/2},-0.95) {#5};}
% box plot centred at height y: \bplotat{y}{min}{Q1}{med}{Q3}{max}
\newcommand{\bplotat}[6]{%
  \draw[thick] (#2,#1) -- (#3,#1);
  \draw[thick] (#5,#1) -- (#6,#1);
  \draw[thick] (#2,#1-0.22) -- (#2,#1+0.22);
  \draw[thick] (#6,#1-0.22) -- (#6,#1+0.22);
  \draw[thick,fill=gray!20] (#3,#1-0.42) rectangle (#5,#1+0.42);
  \draw[very thick] (#4,#1-0.42) -- (#4,#1+0.42);}
% an outlier dot: \outlier{x}{y}
\newcommand{\outlier}[2]{\fill (#1,#2) circle (2.4pt);}

\usepackage[margin=0.75in]{geometry}
\usepackage{array}
\usepackage{enumitem}
\usepackage{amsmath}
\usepackage[helvet]{sfmath}  % math in the same sans font as the text
\usepackage{standalone}
\pagestyle{empty}
\setlength{\parindent}{0pt}
\setlist[enumerate,1]{leftmargin=1.5em,itemsep=1em}
\setlist[enumerate,2]{label=(\alph*),leftmargin=2em,itemsep=2.2em}

\newif\ifkey
\ifdefined\iskey \keytrue \fi
\colorlet{keyred}{red!75!black}
\newcommand{\ans}[1]{\ifkey\leavevmode{\color{keyred}#1}\fi}
% vertical work space on the student copy, less on the key
\newcommand{\work}[1]{\ifkey\vspace{0.6em}\else\vspace{#1}\fi}
\newcommand{\fig}[1]{\input{../../figures/#1}}
% Title line: \lessontitle{Quartiles, IQR, and Box Plots}
\newcommand{\lessontitle}[1]{{\Large\bfseries #1\ifkey\ \textcolor{keyred}{KEY}\fi}\hfill Name \underline{\hspace{2in}}\par\vspace{0.6em}}
\newcommand{\partner}{\par\vspace{0.4em}Work with a partner. Compare answers when you finish each problem.\par}

% Unit quiz: two questions per mini lesson, each headed by the lesson's
% objective.  The key keeps every box the same height as the student copy,
% so key page N lines up with quiz page N.
\newcommand{\lessonhead}[1]{\par\vspace{0.6em}{\large\bfseries #1}\par\vspace{0.25em}\hrule\vspace{0.6em}}
\newcommand{\goal}[1]{\textbf{#1}\par\vspace{0.35em}}
% answer box: same height on both copies, answer printed only on the key
\newcommand{\abox}[2]{\par\begin{minipage}[t][#1][t]{\linewidth}\raggedright\vspace{0.2em}\ans{#2}\end{minipage}\par}
% circled choices: \opt{A} prints plain; \optans{B} is circled in red on the key
\newcommand{\opt}[1]{\tikz[baseline=(c.base)]\node[rounded corners=0.6em,inner sep=2.5pt](c){#1};}
\newcommand{\optans}[1]{\ifkey\tikz[baseline=(c.base)]\node[rounded corners=0.6em,inner sep=2.5pt,draw=keyred,thick,text=keyred](c){#1};\else\opt{#1}\fi}
% blanks on a tree diagram: same width on both copies, filled in red on the key
\newcommand{\tb}[1]{\underline{\makebox[0.42in]{\ans{#1}}}}
\newcommand{\te}[1]{\hspace{0.15in}\underline{\makebox[1.2in][l]{\ans{#1}}}}
\setlist[enumerate,2]{label=(\alph*),leftmargin=2em,itemsep=0.3em}
\setlist[enumerate,1]{leftmargin=1.5em,itemsep=1em}
\raggedright\hyphenpenalty=10000\exhyphenpenalty=10000

\begin{document}

\lessontitle{Unit 6 Quiz: Probability}

Show your work.

% ======================================================================
\lessonhead{6.1 - Venn Diagrams and the Addition Rule}
\begin{enumerate}

\item \goal{I find a probability and its complement, and I read the symbols $A \cup B$ and $A \cap B$.}
40 students said if they have a dog ($D$) and if they have a cat ($C$).

\begin{minipage}[t]{0.44\linewidth}\vspace{0pt}\raggedright
\scalebox{0.72}{\venn{Dog ($D$)}{Cat ($C$)}{13}{5}{10}{12}}
\end{minipage}\hfill
\begin{minipage}[t]{0.53\linewidth}\vspace{0pt}\raggedright
\begin{enumerate}
  \item Find $P(D)$.
  \abox{2.8em}{$\frac{13 + 5}{40} = \frac{18}{40} = \frac{9}{20} = 0.45$}
  \item Find $P(\text{not } D)$.
  \abox{2.8em}{$1 - \frac{9}{20} = \frac{11}{20} = 0.55$}
\end{enumerate}
\end{minipage}

\vspace{0.4em}
\begin{enumerate}[start=3]
  \item What does $D \cap C$ mean in words? Find $P(D \cap C)$.
  \abox{3em}{Has a dog \textbf{and} has a cat. \quad $\frac{5}{40} = \frac{1}{8} = 0.125$}
  \item What does $D \cup C$ mean in words? Find $P(D \cup C)$.
  \abox{3em}{Has a dog \textbf{or} a cat \textbf{or} both. \quad $\frac{13 + 5 + 10}{40} = \frac{28}{40} = \frac{7}{10} = 0.7$}
\end{enumerate}

\item \goal{I use the addition rule, and I decide if two events are mutually exclusive.}
\begin{enumerate}
  \item $P(A) = 0.35$, \; $P(B) = 0.45$, and $P(A \cap B) = 0.15$. Find $P(A \cup B)$.
  \abox{2.6em}{$P(A \cup B) = 0.35 + 0.45 - 0.15 = 0.65$}
\end{enumerate}

\vspace{0.2em}
Draw one card from a deck of 52. Face cards are J, Q, and K. \\
Are the events mutually exclusive? Find the probability.
\begin{enumerate}[start=2]
  \item a 7 or a queen
  \abox{4.3em}{Yes. A card cannot be a 7 and a queen. \\[0.3em] $\frac{4}{52} + \frac{4}{52} = \frac{8}{52} = \frac{2}{13} \approx 0.15$}
  \item a club or a face card
  \abox{4.3em}{No. The J, Q, and K of clubs are both. \\[0.3em] $\frac{13}{52} + \frac{12}{52} - \frac{3}{52} = \frac{22}{52} = \frac{11}{26} \approx 0.42$}
\end{enumerate}

\end{enumerate}

% ======================================================================
\newpage
\lessonhead{6.2 - The Multiplication Rule}
\begin{enumerate}[resume]

\item \goal{I use the multiplication rule for independent events and the complement for ``at least one'' questions.}
A school bus is on time on 90\% of days. The days are independent. Look at 3 school days in a row.
\begin{enumerate}
  \item Find $P(\text{on time all 3 days})$.
  \abox{3.2em}{$0.9 \cdot 0.9 \cdot 0.9 = 0.729$}
  \item Find $P(\text{late all 3 days})$.
  \abox{3.2em}{$0.1 \cdot 0.1 \cdot 0.1 = 0.001$}
  \item Find $P(\text{late at least once})$.
  \abox{3.6em}{$1 - P(\text{on time all 3 days}) = 1 - 0.729 = 0.271$}
\end{enumerate}

\item \goal{I decide if two events are independent or dependent, and I use the general multiplication rule.}
A box has 5 red pens and 7 blue pens. You take 2 pens without looking.
\begin{enumerate}
  \item You put the first pen back before you take the second. Find $P(\text{red, red})$.
  \abox{3.6em}{Independent. \; $\frac{5}{12} \cdot \frac{5}{12} = \frac{25}{144} \approx 0.17$}
  \item You keep the first pen out. Find $P(\text{red second} \mid \text{red first})$.
  \abox{3.6em}{4 red pens are left out of 11 pens. \; $\frac{4}{11} \approx 0.36$}
  \item You keep the first pen out. Find $P(\text{red, red})$.
  \abox{3.6em}{Dependent. \; $\frac{5}{12} \cdot \frac{4}{11} = \frac{20}{132} = \frac{5}{33} \approx 0.15$}
  \item Which pair of events is dependent? Circle the letter.

  \vspace{0.3em}
  \begin{tabular}{@{}ll}
  \opt{A} & Roll a number cube two times.\\
  \opt{B} & Flip a coin, then spin a spinner.\\
  \opt{C} & Pick a name from a hat, put it back, and pick again.\\
  \optans{D} & Pick a name from a hat, keep it out, and pick again.\\
  \end{tabular}
\end{enumerate}

\end{enumerate}

% ======================================================================
\newpage
\lessonhead{6.3 - Conditional Probability}
\begin{enumerate}[resume]

\item \goal{I find conditional probabilities from a two-way table, and I check if two events are independent.}
80 students said if their earbuds are wired or wireless.

\vspace{0.3em}
\begin{minipage}[t]{0.44\linewidth}\vspace{0pt}\raggedright
\renewcommand{\arraystretch}{1.4}
\begin{tabular}{|l|c|c|c|}\hline
 & \textbf{Wired} & \textbf{Wireless} & \textbf{Total} \\ \hline
\textbf{9th grade} & 24 & 16 & 40 \\ \hline
\textbf{10th grade} & 12 & 28 & 40 \\ \hline
\textbf{Total} & 36 & 44 & 80 \\ \hline
\end{tabular}
\end{minipage}\hfill
\begin{minipage}[t]{0.5\linewidth}\vspace{0pt}\raggedright
\begin{enumerate}
  \item Find $P(\text{wireless} \mid \text{10th grade})$.
  \abox{3.8em}{Only the 40 in 10th grade. \\[0.2em] $\frac{28}{40} = \frac{7}{10} = 0.7$}
  \item Find $P(\text{10th grade} \mid \text{wireless})$.
  \abox{3.8em}{Only the 44 with wireless. \\[0.2em] $\frac{28}{44} = \frac{7}{11} \approx 0.64$}
\end{enumerate}
\end{minipage}

\vspace{0.4em}
\begin{enumerate}[start=3]
  \item Are having wireless earbuds and being in 10th grade independent? Show your work.
  \abox{4.4em}{$P(\text{wireless}) = \frac{44}{80} = 0.55$, \; but \; $P(\text{wireless} \mid \text{10th grade}) = 0.7$. \\[0.3em]
  They are different, so the events are \textbf{not} independent.}
\end{enumerate}

\item \goal{I use a tree diagram and the notation $P(A \mid B)$ to find probabilities.}
The chance of rain on Saturday is 30\%. If it rains, the chance that the soccer game is canceled is 80\%.
If it does not rain, the chance that the game is canceled is 5\%.
\begin{enumerate}
  \item Fill in the tree diagram. Write the missing probabilities and multiply along each path.

  \vspace{0.3em}
  \begin{center}\tree{Rain}{0.3}{No rain}{\tb{0.7}}{Canceled}{Played}{0.8}{\tb{0.2}}{0.05}{\tb{0.95}}%
  {\te{$0.3 \cdot 0.8 = 0.24$}}{\te{$0.3 \cdot 0.2 = 0.06$}}{\te{$0.7 \cdot 0.05 = 0.035$}}{\te{$0.7 \cdot 0.95 = 0.665$}}\end{center}
  \item Find $P(\text{canceled})$.
  \abox{3em}{Add the two canceled paths. \; $0.24 + 0.035 = 0.275$}
  \item The game was canceled. Find $P(\text{rain} \mid \text{canceled})$.
  \abox{3.6em}{$\frac{0.24}{0.275} = \frac{48}{55} \approx 0.87$}
\end{enumerate}

\end{enumerate}

% ======================================================================
\newpage
\lessonhead{6.4 - Simulations}
\begin{enumerate}[resume]

\item \goal{I tell theoretical and experimental probability apart, and I explain the law of large numbers.}
A spinner has 5 equal parts. One part is blue. Lena spins it 50 times. It lands on blue 14 times.
\begin{enumerate}
  \item Find the experimental probability of blue.
  \abox{3em}{$\frac{14}{50} = \frac{7}{25} = 0.28$}
  \item Find the theoretical probability of blue.
  \abox{3em}{$\frac{1}{5} = 0.2$}
  \item Lena spins it 5,000 times. What number do you expect the experimental probability of blue to be close to? Why?
  \abox{3.6em}{Close to 0.2. The law of large numbers: with more trials, the experimental probability gets closer to the theoretical probability.}
\end{enumerate}

\item \goal{I assign random digits and run a simulation to estimate a probability.}
A soccer team wins 60\% of its games. Simulate the next 4 games. \\
Digits 0 to 5 = win. \; Digits 6 to 9 = loss. \; One trial = 4 digits (4 games).
\begin{enumerate}
  \item Why do the digits 0 to 5 stand for a win?
  \abox{2.6em}{0 to 5 is 6 of the 10 digits, and $\frac{6}{10} = 60\%$.}
  \item Use these random digits to run 10 trials. Count the wins in each trial.

  \begin{center}
  {\large\texttt{\bfseries 76869 45827 19553 09040 57143 14173 96580 44469}}

  \vspace{0.6em}
  \renewcommand{\arraystretch}{1.6}
  \setlength{\tabcolsep}{4pt}
  \begin{tabular}{|l|p{0.42in}|p{0.42in}|p{0.42in}|p{0.42in}|p{0.42in}|p{0.42in}|p{0.42in}|p{0.42in}|p{0.42in}|p{0.42in}|}\hline
  Trial & 1 & 2 & 3 & 4 & 5 & 6 & 7 & 8 & 9 & 10 \\ \hline
  Digits & \ans{7686} & \ans{9458} & \ans{2719} & \ans{5530} & \ans{9040} & \ans{5714} & \ans{3141} & \ans{7396} & \ans{5804} & \ans{4469} \\ \hline
  Wins & \ans{0} & \ans{2} & \ans{2} & \ans{4} & \ans{3} & \ans{3} & \ans{4} & \ans{1} & \ans{3} & \ans{2} \\ \hline
  \end{tabular}
  \end{center}
  \item Use your simulation to estimate $P(\text{at least 3 wins})$.
  \abox{3em}{Trials 4, 5, 6, 7, and 9 have 3 or more wins. \; $\frac{5}{10} = 0.5$}
  \item A different team wins 35\% of its games. How would you assign random digits to one game?
  \abox{3em}{Use pairs of digits. \; 00 to 34 = win (35 of the 100 pairs). \; 35 to 99 = loss.}
\end{enumerate}

\end{enumerate}

% ======================================================================
\newpage
\lessonhead{6.5 - Permutations}
\begin{enumerate}[resume]

\item \goal{I use the counting principle and factorials.}
\begin{enumerate}
  \item A bike comes in 4 colors, 3 sizes, and 2 seat styles. How many different bikes are there?
  \abox{3.6em}{$4 \cdot 3 \cdot 2 = 24$ bikes}
  \item Find $\dfrac{7!}{5!}$.
  \abox{3.6em}{$\frac{7 \cdot 6 \cdot 5 \cdot 4 \cdot 3 \cdot 2 \cdot 1}{5 \cdot 4 \cdot 3 \cdot 2 \cdot 1} = 7 \cdot 6 = 42$}
  \item In how many ways can 7 songs be put in order on a playlist?
  \abox{3.6em}{$7! = 7 \cdot 6 \cdot 5 \cdot 4 \cdot 3 \cdot 2 \cdot 1 = 5{,}040$ ways}
\end{enumerate}

\item \goal{I count permutations with ${}_nP_r$, and I tell when repeats are allowed.}
\begin{enumerate}
  \item 9 students run for president, vice president, and secretary. In how many ways can the 3 jobs be filled?
  \abox{3.6em}{Order matters. \; ${}_9P_3 = \frac{9!}{6!} = 9 \cdot 8 \cdot 7 = 504$ ways}
\end{enumerate}
A locker code uses 4 letters from A to F (6 letters).
\begin{enumerate}[start=2]
  \item How many codes are there if letters can repeat?
  \abox{4em}{$6 \cdot 6 \cdot 6 \cdot 6 = 6^4 = 1{,}296$ codes}
  \item How many codes are there if no letter repeats?
  \abox{4em}{${}_6P_4 = 6 \cdot 5 \cdot 4 \cdot 3 = 360$ codes}
\end{enumerate}

\end{enumerate}

% ======================================================================
\newpage
\lessonhead{6.6 - Combinations}
\begin{enumerate}[resume]

\item \goal{I count combinations with ${}_nC_r$.}
\begin{enumerate}
  \item Find ${}_6C_4$.
  \abox{3.6em}{$\frac{6!}{4!\,2!} = \frac{6 \cdot 5}{2 \cdot 1} = 15$}
  \item 11 teams are in a league. Each team plays every other team once. How many games are there?
  \abox{3.6em}{A game is a pair of teams. \; ${}_{11}C_2 = \frac{11 \cdot 10}{2} = 55$ games}
  \item You choose 4 of 9 photos for a poster. In how many ways can you choose them?
  \abox{3.6em}{${}_9C_4 = \frac{9 \cdot 8 \cdot 7 \cdot 6}{4 \cdot 3 \cdot 2 \cdot 1} = 126$ ways}
\end{enumerate}

\item \goal{I decide whether a problem needs a permutation or a combination, and I explain the difference.}
For each situation, write permutation or combination. Then find the number of ways.

\vspace{0.4em}
{\renewcommand{\arraystretch}{1.8}
\begin{tabular}{|>{\raggedright\arraybackslash}p{2.9in}|>{\raggedright\arraybackslash}p{1.3in}|>{\raggedright\arraybackslash}p{1.7in}|}\hline
\textbf{Situation} & \textbf{Permutation or combination?} & \textbf{Number of ways} \\ \hline
(a) Choose 3 of 11 students for a clean-up crew. & \ans{combination} & \ans{${}_{11}C_3 = 165$} \\ \hline
(b) Give 1st, 2nd, and 3rd place to 3 of 11 science projects. & \ans{permutation} & \ans{${}_{11}P_3 = 990$} \\ \hline
\end{tabular}}
\begin{enumerate}[start=3]
  \item Why is the answer to (b) 6 times the answer to (a)?
  \abox{3.6em}{Each group of 3 can be put in order $3! = 6$ ways. A permutation counts every order. A combination counts the group once.}
  \item Which situation is a combination? Circle the letter.

  \vspace{0.3em}
  \begin{tabular}{@{}ll}
  \opt{A} & Choosing a captain and a co-captain from 15 players.\\
  \optans{B} & Choosing 5 of 15 players to ride in the first van.\\
  \opt{C} & Arranging 5 trophies in a row on a shelf.\\
  \opt{D} & Making a 4-digit code with no repeated digits.\\
  \end{tabular}
\end{enumerate}

\end{enumerate}

% ======================================================================
\newpage
\lessonhead{6.7 - Probability with Counting}
\begin{enumerate}[resume]

\item \goal{I count the possible outcomes and the outcomes I want, and I find the probability.}
A coach picks 3 players at random from 7 forwards and 5 defenders.
\begin{enumerate}
  \item How many groups of 3 are possible?
  \abox{3.6em}{${}_{12}C_3 = \frac{12 \cdot 11 \cdot 10}{3 \cdot 2 \cdot 1} = 220$ groups}
  \item Find $P(\text{all 3 are forwards})$.
  \abox{4.4em}{${}_7C_3 = 35$ groups of forwards. \; $\frac{35}{220} = \frac{7}{44} \approx 0.16$}
  \item Find $P(\text{all 3 are defenders})$.
  \abox{4.4em}{${}_5C_3 = 10$ groups of defenders. \; $\frac{10}{220} = \frac{1}{22} \approx 0.05$}
\end{enumerate}

\item \goal{I decide if order matters in a probability problem.}
15 students are in a club. Does order matter? Find the probability.
\begin{enumerate}
  \item 3 students are chosen at random to be president, vice president, and secretary. What is the probability that Amira is president, Ben is vice president, and Carla is secretary?
  \abox{5.4em}{Yes, the jobs are different. \; ${}_{15}P_3 = 2{,}730$. \; $P = \frac{1}{2{,}730} \approx 0.0004$}
  \item 3 students are chosen at random for a committee. What is the probability that the committee is Amira, Ben, and Carla?
  \abox{5.4em}{No, it is the same job. \; ${}_{15}C_3 = 455$. \; $P = \frac{1}{455} \approx 0.002$}
\end{enumerate}

\end{enumerate}

\end{document}
